\documentclass[11pt]{article}
\usepackage{fontspec,amsmath,amssymb,amsthm,hyperref}
\usepackage[margin=1in]{geometry}
\setlength{\emergencystretch}{3em}
\title{Normalization, moments and score moments of the lifted target}
\author{LangevinLean mathematical exposition}
\date{}
\begin{document}\maketitle
\noindent Source module: \texttt{IPLA/Foundations.lean}. Compile this proof document with XeLaTeX or LuaLaTeX. The accompanying HTML places the complete Lean declarations beside every proof.


Apply the shared ULA Gibbs results to \(V_N\), whose differentiability and \(Nm\)-strong convexity were derived from \(U\). All normalization and moment facts are conclusions from the original joint assumptions.


\section{The named target is the normalized Gibbs measure}


The definitions give \(\operatorname{gibbsTarget}(U)=\operatorname{gibbsMeasure}(V_N)\). If \(N>0\) and \(m>0\), then \(Nm>0\). The lifted strong-convexity theorem and chain-rule differentiability therefore meet the shared Gibbs theorem's hypotheses. Its quadratic lower bound gives integrable Gaussian tails and a positive partition, proving \(0<Z_N<\infty\) and thus \(\Pi_N(\mathbb R^D)=1\). The same tails imply \(\int\|z\|^2d\Pi_N<\infty\), expressed as membership of the identity in \(L^2(\Pi_N)\).



\noindent Lean declarations: \texttt{gibbsTarget\_eq\_gibbsMeasure}, \texttt{gibbsTarget\_isProbability}, \texttt{gibbsTarget\_memLp}.

\section{Finite transport supplies square integrability of initialization}


Suppose a measurable initialization \(X\) has \(W_2(X_\#P,\Pi_N)<\infty\). The target already has a finite second moment. For any finite-cost coupling, \(\|x\|^2\le2\|x-y\|^2+2\|y\|^2\); integrating transfers the target moment and coupling cost to the initial law. The shared transport-moment theorem formalizes this argument and proves \(X\in L^2(P)\). Thus this alternative initialization interface introduces no independent moment assumption.



\noindent Lean declarations: \texttt{initial\_memLp\_of\_finite\_transport}.

\section{The lifted Gibbs score moment has the correct dimension}


The shared score estimate for a strongly convex potential with \(\ell\)-Lipschitz gradient is \(\int\|\nabla V\|^2d\pi\le\ell\dim H\). Substitute \(V=V_N\), \(\ell=NL\), and \(\dim H=d_\theta+Nd_x\) to obtain

\[\int\|\nabla V_N(z)\|^2d\Pi_N(z)\le NL(d_\theta+Nd_x).\]

The lifted smoothness and convexity hypotheses follow from the joint ones. The proof only rewrites the common Gibbs definition and exact dimension; it adds no score-moment or Hessian premise.



\noindent Lean declarations: \texttt{gibbsTarget\_score\_moment}.

\end{document}
